\section{Claim Audit: What the Record Will Bear}

\begin{longtable}{>{\bfseries\RaggedRight\arraybackslash}p{0.25\linewidth}>{\RaggedRight\arraybackslash}p{0.19\linewidth}>{\RaggedRight\arraybackslash}p{0.47\linewidth}}
\toprule
\textbf{Shorthand claim} & \textbf{Verdict} & \textbf{Source-controlled formulation}\\
\midrule
\endfirsthead
\toprule
\textbf{Shorthand claim} & \textbf{Verdict} & \textbf{Source-controlled formulation}\\
\midrule
\endhead
\bottomrule
\endfoot
“The SSPX was illegal from the beginning.” & Rejected. & Bishop Charrière erected a diocesan \emph{pia unio} in 1970 and approved its statutes for six years experimentally. The later question is how that bounded approval ended and whether subsequent ministry had a canonical basis.\\
“Rome permanently approved it in 1971.” & Rejected. & Cardinal Wright praised the statutes in a letter. No checked text makes that letter a decree of pontifical-right erection or a permanent universal mission.\\
“The 1975 suppression is an uncontested legal fact.” & Overstated. & The diocesan and Roman authorities treated it as effective and Paul VI confirmed that position; Lefebvre and the SSPX contested competence and procedure. No checked competent judgment restored the status.\\
“The conflict was only about preserving the Latin Mass.” & Rejected. & Liturgy was central, but contemporary papal and SSPX documents also concern conciliar doctrine, magisterial authority, seminary governance, ordination, and ecclesial obedience.\\
“Lefebvre rejected the pope and became sedevacantist.” & Rejected. & He continued to claim recognition of the Roman Pontiff. In 1983 SSPX leadership expelled priests whose positions it described as tending toward denial of ordinary papal legitimacy and enforced a boundary against those conclusions. That fact does not resolve the charge of schism.\\
“Rome refused any SSPX bishop in 1988.” & Rejected. & The 5 May protocol contemplated episcopal provision through the papal appointment process. The accord collapsed over confidence, timing, guarantees, and Lefebvre's decision to proceed without mandate.\\
“The 1988 consecrations were invalid.” & Rejected as stated. & The official dispute concerned illicit consecration without mandate, schism, and penalties. Absence of papal mandate is not presented in the checked acts as absence of sacramental episcopal orders.\\
“Every SSPX priest and attendee was excommunicated in 1988.” & Rejected. & The decree named the consecrating and consecrated bishops. The 1996 note required internal and external formal adherence and case-specific assessment; occasional attendance alone was insufficient.\\
“The 2009 remission regularized the SSPX.” & Rejected. & It remitted the four living bishops' personal censures. Benedict XVI expressly said the Society had no canonical status and its ministers did not legitimately exercise ministry.\\
“Francis later gave the SSPX ordinary faculties.” & Misleading. & Francis supplied a confession faculty for a defined pastoral purpose and enabled local ordinaries to delegate marriage assistance through a structured route. These were real but bounded concessions, not canonical erection.\\
“All SSPX sacraments are invalid.” & Rejected. & The 2026 note calls SSPX sacramental celebrations illicit and identifies penance and marriages assisted by their ministers as invalid. Validity is sacrament-specific.\\
“Every person attending an SSPX Mass after July 2026 is excommunicated.” & Rejected. & The DDF retained the 1996 formal-adherence test for lay persons. Intention, external manifestation, and individual facts matter; occasional attendance is not automatically formal adherence.\\
“The SSPX appeal has already erased the 2 July acts.” & Unsupported. & The Society announced preliminary recourse against the penal decree and claimed suspensive effect. No Holy See response had been located by 16 July, and the communiqué did not establish that the separate explanatory note was suspended.\\
\end{longtable}

\subsection{Findings by confidence}

\begin{longtable}{>{\bfseries\RaggedRight\arraybackslash}p{0.36\linewidth}>{\RaggedRight\arraybackslash}p{0.16\linewidth}>{\RaggedRight\arraybackslash}p{0.39\linewidth}}
\toprule
\textbf{Finding} & \textbf{Confidence} & \textbf{Reason and limit}\\
\midrule
\endhead
The Society received genuine diocesan erection as a \emph{pia unio} in 1970, while its statutes received bounded experimental approval. & High. & The dated act distinguishes erection from the six-year approval, tacit renewal, and possible later definitive erection; its exact legal consequences under the 1917 Code still require specialist treatment beyond this history.\\
The 1974--1976 conflict joined doctrine, liturgy, and governance. & High. & Participant declarations, suspensions, and Paul VI's acts converge; motives of every individual participant do not.\\
The 1988 protocol offered a substantial canonical structure and episcopal provision. & High. & The signed text survives; counterfactual claims about whether it would have succeeded remain unknowable.\\
The 1988 consecrations secured institutional episcopal continuity while deepening rupture. & High. & The act, named participants, penalties, subsequent ordinations, and new communities are documented.\\
The 2009--2012 process failed principally over doctrinal agreement, not lack of an imaginable structure. & High. & Benedict's acts and CDF communiqués explicitly frame the problem; private negotiating motives remain only partially public.\\
Francis's faculties relieved specific sacramental problems without regularization. & High. & The papal letter, apostolic letter, and marriage instruction define both grants and limits.\\
The DDF made an express prospective determination concerning SSPX sacred ministers and current sacramental consequences on 2 July 2026. & High as an official-act report. & The Italian note is explicit and incorporates the 1996 analysis; its application to individual cases and the effect of pending recourse demand competent canonical judgment.\\
SSPX growth proves its theological claims. & Not established. & Self-reported scale proves institutional persistence, not divine approval, doctrinal correctness, or canonical mission.\\
\bottomrule
\end{longtable}
